\scanpage{76}
% Archive pencil foliation: 75.

part, \inserted{\handwritten{\uline{une fois choisi} \uline{une clôture} $\overline{\mathbb{F}}_p$,}} la gerbe $\handwritten{\underline{\pi}}_p(\underline{F}_p)$, grâce àu\deleted{son} foncteur \unsure{\deleted{e}f|\deleted{f}e|f}\unsure{\deleted{v}i|\deleted{i}v|v|i|ï}bre évident ``fibre en $\handwritten{\bar{\underline{F}}}_p$ '' sur la catégorie des $\underline{Q}_p$-faisceaux constructibles sur $\underline{F}_p$, s'identifie à la \deleted{catégorie}\inserted{gerbe triviale défi\unsure{\deleted{m}n|\deleted{n}m|n|m}\unsure{\deleted{t}i|\deleted{i}t|\deleted{e}i|\deleted{i}e|i|t}\unsure{\deleted{a}e|\deleted{e}a|e|a}\deleted{\unsure{et|alt|elt|ilt}} par .l'}enveloppe $\underline{Q}_p$-alg \unsure{e|é}briqueque du gro upe profini $\handwritten{\pi}_1(\underline{F}_p)=\handwritten{\bar{\underline{Z}}}\deleted{\unsure{r|n|\bar r}}$, \deleted{lequel, comme on le constate aisément, est le}

Un \inserted{$\underline{Q}_p$-}faisceau constructible sur $\underline{F}_p$ s'identifie à un vectoriel de dimension finie $E$ sur $\underline{Q}_p$, muni d'un automorphisme $F_E$ tel que l'\handwritten{\unsure{a|o}}\deleted{\unsure{hom|lun}} opération $n\handwritten{\mapsto}F_E^n$ de $\underline{Z}$ sur $E$ se prolonge par \unsure{co|\deleted{c}o|\deleted{o}c|$\varphi$} ntinuité en une opération de \deleted{\unsure{$\hat{\underline{Z}}_{\unsure{p|r|E}}$|$\hat{\underline{Q}}_{\unsure{p|r|E}}$|$\hat{\underline{Z}}r$}} $\handwritten{\hat{\underline{Z}}}$ \margin{\handwritten{$Z$}}. Identifiant \inserted{\handwritten{(par un foncteur $j$, \uline{distinct} de celui de 5\unsure{,|.}2)}} \unsure{\deleted{d}l|\deleted{l}d|l|d}a catégorie de ces $(E,\unsure{\mathbb{E}_E|\deleted{E_E}|\deleted{F}E_E|\deleted{E}F_E}F_E)$ à une\handwritten{/}sous-catégorie pleine de\deleted{l} la catégorie \deleted{des} $\mathrm{Fcriso}(\underline{Q}_p)$ des couples $(E,F_E)$, où $F_E$ est un aut omorphisme \uline{quelconque}, on voit que le lien \deleted{de notre}\inserted{$G\handwritten{'}$ de} $\handwritten{\underline{\pi}}_p(\underline{F}_p)$ de cette catégorie s'identifie à un quotient du lien $G$ de $\mathrm{Fcriso}(\underline{F}_p)$ déterminé plus haut, et on voit aussitôt que ce quotient est de la forme $G_u\mathrm{x}G_r^{\overset{\handwritten{\unsure{2|\ell}}\deleted{\unsure{\prime|1|x}}}{\handwritten{'}}}$, où $G_r'$ est \unsure{\deleted{u}l|\deleted{l}u|\deleted{h}l|\deleted{l}h|l|u|h}\unsure{\deleted{a}e|\deleted{e}a|e|a|n} qutient de $G_r$ (ceci revient à dire que si on décompose $F_E$ en produit de sa composante unipotente par sa composante semi-simple, alors la condition que $F_E$ \unsure{\deleted{e}d|\deleted{d}e|\deleted{s}d|\deleted{d}s|d}éf\unsure{\deleted{t}i|\deleted{i}t|\deleted{r}i|\deleted{i}r|i}nisse une opération de $\handwritten{\hat{\underline{Z}}}$ ne dépend que de la composante semi-simple); de façon précise, $G'_r$ est le groupe quotient de $G_r$ défini par le sous-groupe du groupe des caractères $\handwritten{\bar{\underline{Q}}}_p^{\handwritten{*}}$ de $G_r$ formé des \uline{unités} p\unsure{\deleted{é}-|\deleted{-}é|é|-}\uline{adiques} : cela signifie en effe\unsure{\deleted{y}t|\deleted{t}y|t} que $F_E$ définit uneopération de \deleted{$\bar{\unsure{F|G}}_E$} $\handwritten{\hat{\underline{Z}}}$ (i.e. définit un F-isocristal de pente nulle) sss les valeur\unsure{\deleted{e}s|\deleted{s}e|s|e}\deleted{\unsure{s|r}} propres de $F_E$ sont des unités, ce qui ré\unsure{su|us}lte en effet aussitôt du critère donné plus haut des F-isocristaux de pente nulle (ou se vérifie directement de façon immédiate; la référence au résultat sur les fourbis de pente nulle devrait se justifier par un grain de sel, car l'inclusion \unsure{\deleted{y}u|\deleted{u}y|u}tilisée ici des $\underline{Q}_p$-faisceaux constructibles sur $\underline{F}_p$ dans\deleted{l} les F-iso\deleted{cr}cristaux n'est pas celle utilisée \deleted{plus h\unsure{\deleted{u}a|\deleted{a}u|a|u}\unsure{\deleted{y}u|\deleted{u}y|u}t}\inserted{\handwritten{dans 5.2}}, cf. plus bas). On trouve ain\unsure{\deleted{i}s|\deleted{s}i|s}\unsure{\deleted{s}i|\deleted{i}s|i} une structure d'extension sur $G\handwritten{\unsure{=|\simeq|\ne}}\operatorname{Lien}(\underline{G}(\underline{F}_p))$:
\diagram{\begin{tikzcd}[column sep=large,execute at end picture={\draw[gray,thin] ([xshift=5pt,yshift=3pt]\tikzcdmatrixname-1-4.east)--([xshift=10pt,yshift=-1pt]\tikzcdmatrixname-1-4.east);}]0\arrow[r]&H\arrow[r]&G\arrow[r]&G'\arrow[r]&0\ ,\end{tikzcd}}
où $H$ est le groupe de t.m. dont le groupe des caractères est le
